\par\Needspace{6\baselineskip}\medskip
\hypertarget{native-130-article096}{}\pdfbookmark[2]{96.}{article-096}
\begin{center}96.\end{center}
Aequationes in art. praec. traditae tanta quidem concinnitate gaudent, ut nihil amplius desiderari posse videatur. Nihilominus eruere licet formulas quasdam alias, per quas elementa orbitae multo adhuc elegantius et commodius determinantur: verum evolutio harum formularum paullulo magis recondita est.

Resumimus ex art. 8 aequationes sequentes, quas commoditatis gratia numeris novis distinguimus:
\[
\mathrm{I.}\qquad \sin\tfrac12 v\mathpunct{.}\BareRoot{\dfrac{r}{a}}=\sin\tfrac12 E\mathpunct{.}\BareRoot{(1+e)}
\]
\SourceLeaf{131}
\[
\begin{aligned}
\mathrm{II.}\qquad \cos\tfrac12 v\mathpunct{.}\BareRoot{\dfrac{r}{a}}&=\cos\tfrac12 E\mathpunct{.}\BareRoot{(1-e)}\\[6pt]
\mathrm{III.}\qquad \sin\tfrac12 v'\mathpunct{.}\BareRoot{\dfrac{r'}{a}}&=\sin\tfrac12 E'\mathpunct{.}\BareRoot{(1+e)}\\[6pt]
\mathrm{IV.}\qquad \cos\tfrac12 v'\mathpunct{.}\BareRoot{\dfrac{r'}{a}}&=\cos\tfrac12 E'\mathpunct{.}\BareRoot{(1-e)}.
\end{aligned}
\]
Multiplicamus I per $\sin\frac12(F+g)$, II per $\cos\frac12(F+g)$, unde productis additis nanciscimur
\[
\begin{aligned}
\cos\tfrac12(f+g)\mathpunct{.}\BareRoot{\dfrac{r}{a}}&=\sin\tfrac12 E\sin\tfrac12(F+g)\mathpunct{.}\BareRoot{(1+e)}\\[4pt]
&\quad+\cos\tfrac12 E\cos\tfrac12(F+g)\mathpunct{.}\BareRoot{(1-e)}
\end{aligned}
\]
sive, propter $\BareRoot{(1+e)}=\cos\frac12\varphi+\sin\frac12\varphi$, $\BareRoot{(1-e)}=\cos\frac12\varphi-\sin\frac12\varphi$:
\[
\begin{aligned}
\cos\tfrac12(f+g)\mathpunct{.}\BareRoot{\dfrac{r}{a}}&=\cos\tfrac12\varphi\cos(\tfrac12 F-\tfrac12 G+g)\\[4pt]
&\quad-\sin\tfrac12\varphi\cos\tfrac12(F+G).
\end{aligned}
\]
Prorsus simili modo multiplicando III per $\sin\frac12(F-g)$, IV per $\cos\frac12(F-g)$, prodit productis additis
\[
\begin{aligned}
\cos\tfrac12(f+g)\mathpunct{.}\BareRoot{\dfrac{r'}{a}}&=\cos\tfrac12\varphi\cos(\tfrac12 F-\tfrac12 G-g)\\[4pt]
&\quad-\sin\tfrac12\varphi\cos\tfrac12(F+G).
\end{aligned}
\]
Subtrahendo ab hac aequatione praecedentem, oritur
\[
\cos\tfrac12(f+g)\mathpunct{.}\left(\BareRoot{\dfrac{r'}{a}}-\BareRoot{\dfrac{r}{a}}\right)=2\cos\tfrac12\varphi\sin g\sin\tfrac12(F-G)
\]
sive introducendo angulum auxiliarem $\omega$
\[
\begin{aligned}
\mathrm{27)}\qquad &\cos\tfrac12(f+g)\operatorname{tang}2\omega\\[4pt]
&=\sin\tfrac12(F-G)\cos\tfrac12\varphi\sin g\mathpunct{.}\BareRoot[4]{\dfrac{aa}{rr'}}\mathpunct{\cdot}
\end{aligned}
\]
Per transformationes prorsus similes, quarum evolutionem lectori perito relinquimus, invenitur
\[
\begin{aligned}
\mathrm{28)}\qquad &\dfrac{\sin\frac12(f+g)}{\cos 2\omega}\\[4pt]
&=\cos\tfrac12(F-G)\cos\tfrac12\varphi\sin g\mathpunct{.}\BareRoot[4]{\dfrac{aa}{rr'}}\\[8pt]
\mathrm{29)}\qquad &\cos\tfrac12(f-g)\operatorname{tang}2\omega\\[4pt]
&=\sin\tfrac12(F+G)\sin\tfrac12\varphi\sin g\mathpunct{.}\BareRoot[4]{\dfrac{aa}{rr'}}\\[8pt]
\mathrm{30)}\qquad &\dfrac{\sin\frac12(f-g)}{\cos 2\omega}\\[4pt]
&=\cos\tfrac12(F+G)\sin\tfrac12\varphi\sin g\mathpunct{.}\BareRoot[4]{\dfrac{aa}{rr'}}\mathpunct{\cdot}
\end{aligned}
\]
Quum partes primae in his quatuor aequationibus sint quantitates cognitae, ex 27 et 28 determinabuntur $\frac12(F-G)$ et $\cos\frac12\varphi\sin g\mathpunct{.}\BareRoot[4]{\dfrac{aa}{rr'}}=P$, nec non ex 29 et 30 perinde $\frac12(F+G)$ et $\sin\frac12\varphi\sin g\mathpunct{.}\BareRoot[4]{\dfrac{aa}{rr'}}=Q$; ambiguitas in \SourceLeaf{132}determinatione angulorum $\frac12(F-G)$, $\frac12(F+G)$ ita decidenda est, ut $P$ et $Q$ cum $\sin g$ idem signum obtineant. Dein ex $P$ et $Q$ derivabuntur $\frac12\varphi$ et $\sin g\mathpunct{.}\BareRoot[4]{\dfrac{aa}{rr'}}=R$. Ex $R$ deduci potest $a=\dfrac{RR\BareRoot{rr'}}{\sin g^2}\mathpunct{\raisebox{.45ex}{,}}$ nec non $p=\dfrac{\sin f^2\BareRoot{rr'}}{RR}\mathpunct{\raisebox{.45ex}{,}}$ nisi illa quantitate, quae fieri debet $=\pm\BareRoot{\{2(l+\sin\frac12 g^2)\cos f\}}$, $=\pm\BareRoot{\{-2(L-\sin\frac12 g^2)\cos f\}}$, unice ad calculi confirmationem uti malimus, in quo casu $a$ et $p$ commodissime determinantur per formulas
\[
b=\dfrac{\sin f\mathpunct{.}\BareRoot{rr'}}{\sin g}\mathpunct{\raisebox{.45ex}{,}}\qquad a=\dfrac{b}{\cos\varphi}\mathpunct{\raisebox{.45ex}{,}}\qquad p=b\cos\varphi.
\]
Possunt etiam, pro lubito, plures aequationum artt. 88 et 95 ad calculi confirmationem in usum vocari, quibus sequentes adhuc adiicimus:
\[
\begin{aligned}
\dfrac{2\operatorname{tang}2\omega}{\cos 2\omega}\BareRoot{\dfrac{rr'}{aa}}&=e\sin G\sin g\\[6pt]
\dfrac{2\operatorname{tang}2\omega}{\cos 2\omega}\BareRoot{\dfrac{pp}{rr'}}&=e\sin F\sin f\\[6pt]
\dfrac{2\operatorname{tang}2\omega}{\cos 2\omega}&=\operatorname{tang}\varphi\sin G\sin f=\operatorname{tang}\varphi\sin F\sin g.
\end{aligned}
\]
Denique motus medius atque epocha anomaliae mediae perinde invenientur ut in art. praec.
